% ---
% id: protocol-compiler/rem:party-label-restriction
% kind: rem
% title: "Party labels in the general format"
% summary: >-
%   Deriving the party labels from the factor list is narrower than the basic
%   format, which may label the tensor factors of one sector differently. Only
%   the premises placing the laboratories read the labels; two parties inside
%   one measured system need two measured factors.
% prominence: standard
% status: accepted
% uses: [def:declaration]
% mentions: [meat-core/ass:lab, transmitter/ass:src]
% assumes: []
% refs: []
% provenance: newly-derived
% lean:
% ---

\begin{remark}\label{rem:party-label-restriction}
The party labels that (DX1) of Definition~\ref{def:declaration} derives from
the factor list are narrower than the party partition of the basic format at
one measured factor: there a declared tensor factorization of $\Hs{\Qsys_s}$
may carry different party labels on different tensor factors, while in the
general format every tensor factor of one measured factor's block carries
that factor's party. The narrowing is admitted and not repaired. No
derivation of the compiler or of the consuming engine reads those labels ---
what the derivations consume is the readout maps $\rd{\mathsf A}$,
$\rd{\mathsf B}$ --- and correspondingly no check compares a declared label
with a derived one. The labels are read by the premises that place the
parties' laboratories: the consuming engine's laboratory premise
(\xref{meat-core/ass:lab}) takes them, with the party label of $A$, as naming
the parties' laboratories, and the source licence
(\xref{transmitter/ass:src}) reads the party label of $A$; so the narrowing restricts which laboratory boundaries a
general-format declaration can state, and a declaration that wants two parties inside one
\tref{measured system} declares two measured factors instead. So the general format
coincides with the basic format entry for entry on every component a
derivation reads, and on the party labels it is a restriction.
\end{remark}
